\usepackage[utf8]{inputenc}%(only for the pdftex engine)
% \usepackage[big]{dgruyter}
\usepackage{PRIMEarxiv}
 \makeatletter
\newcommand\schapter[1]{%
    \if@openright\cleardoublepage\else\clearpage\fi
    \phantomsection % Ensures hyperref creates a link
    \@schapter{#1}%
    \addcontentsline{toc}{chapter}{#1} % Adds the schapter to the TOC as a chapter
}
\providecommand*\toclevel@schapter{0} % Set as a top-level (like chapter)
\providecommand*\Hy@toclevel@schapter{0} % Ensure hyperref recognizes it at the same level
\makeatother

% \usepackage[nonatbib]{neurips_2023}
% LINE numbers
% \usepackage[modulo]{lineno}
% \linenumbers
% \modulolinenumbers[1]

\usepackage[utf8]{inputenc} % allow utf-8 input
\usepackage[T1]{fontenc}    % use 8-bit T1 fonts
\usepackage{hyperref}       % hyperlinks
\usepackage{url}            % simple URL typesetting
\usepackage{booktabs}       % professional-quality tables
\usepackage{amsfonts}       % blackboard math symbols
\usepackage{nicefrac}       % compact symbols for 1/2, etc.
\usepackage{microtype}      % microtypography
\usepackage{xcolor}         % colors
\usepackage{amsmath}
\usepackage{amsthm}
\usepackage{amssymb}
\usepackage{xparse}
\usepackage{graphicx}
\usepackage{cancel}
\usepackage{comment}
\usepackage{pifont}% http://ctan.org/pkg/pifont
\usepackage{mathtools}
\usepackage[noabbrev,nameinlink]{cleveref}
\usepackage{enumitem}
\setlist[enumerate]{leftmargin=2em}
\setlist[itemize]{leftmargin=2em}
\setlist[description]{leftmargin=*}
\usepackage{tcolorbox}

\RequirePackage{fix-cm}  % Allows more flexible sizes for Computer Modern
\usepackage{anyfontsize}  % Allows arbitrary font sizes

\usepackage{thm-restate} % Allows restating theorems

\hypersetup{
    breaklinks=true,
    pdfencoding=unicode,
    bookmarksnumbered,
    pdfborder={0 0 0},
}%

\newcommand{\cmark}{\ding{51}}%
\newcommand{\xmark}{\ding{55}}%
\definecolor{lightblue}{rgb}{0.6 , 0.7, 0.9}
\definecolor{lightpurple}{rgb}{0.65 , 0.6, 0.9}
\definecolor{lightcyan}{rgb}{0.5 , 0.85, 0.8}
\definecolor{darkpurple}{rgb}{0.35 , 0.15, 0.6}

% Define the theorem-like environment
\newtheorem{theorem}{Theorem}[section]
\newtheorem{definition}[theorem]{Definition}
\newtheorem{example}[theorem]{Example}
\newtheorem{lemma}[theorem]{Lemma}
\newtheorem{corollary}[theorem]{Corollary}
\newtheorem{remark}[theorem]{Remark}
\newtheorem{summary}[theorem]{Summary}
\newtheorem{notation}[theorem]{Notation}
\newtheorem{proposition}[theorem]{Proposition}
% \newtheorem{conjecture}[theorem]{Conjecture}
% \newtheorem{observation}[theorem]{Observation}

\NewDocumentCommand{\optionalbracket}{m m}{%
    \NewDocumentCommand{#1}{g}{%
        \IfNoValueTF{##1}{%
            #2%
        }{%
            #2(##1)%
        }%
    }%
}


\NewDocumentCommand{\oneortwo}{m m m}{%
    \NewDocumentCommand{#1}{g}{%
        \IfNoValueTF{##2}{%
            #3(##1)%
        }{%
            #3(##1,##2)%
        }%
    }%
}

\newcommand{\distrsym}{\bigtriangleup}
\optionalbracket{\distributions}{\distrsym}
\optionalbracket{\distributionsO}{\distrsym_0}
\optionalbracket{\distributionsF}{\distrsym^\otimes}
\optionalbracket{\distributionsFO}{\distrsym^*_0}
\newcommand{\distributionsFC}[2]{\distrsym^*_{#2}(#1)}
\newcommand{\distributionsFi}[2]{\distrsym^*_{#2}(#1)}
\newcommand{\distributionsFCi}[3]{\distrsym^*_{#2,#3}(#1)}
\newcommand{\distributionsFOi}[2]{\distrsym^*_{0,#2}(#1)}


\newcommand{\differenton}[1]{\tildenot^(#1)}
\newcommand{\factors}{\Omega}
\newcommand{\indep}{\mathrel{\perp\!\!\!\perp}}
\newcommand{\orth}{\mathrel{\perp}}
\newcommand{\orthF}{\mathrel{\perp}^\FS}
\newcommand{\orthnot}{\not\orth}%{\cancel{\orth}}
\newcommand{\orthnotF}{\not\orth^\FS}%{\cancel{\orth}}
\newcommand{\sect}{\cap}
\newcommand{\sectbig}{\bigcap}
\newcommand{\union}{\cup}
\newcommand{\unionbig}{\bigcup}
\newcommand{\indepstar}{\indep^\otimes}
\newcommand{\timesbig}{\mathop{\text{\Large$\times$}}}%
\newcommand{\RR}{\mathbb{R}}
\newcommand{\NN}{\mathbb{N}}
\newcommand{\PP}{\mathbb{P}}
\newcommand{\comp}[1]{{\overline{#1}}}
\newcommand{\abs}[1]{|#1|}
\newcommand{\ifff}{\Leftrightarrow}
\newcommand{\impl}{\Rightarrow}
\newcommand{\implied}{\Leftarrow}
\newcommand{\nothing}{\varnothing}
\newcommand{\given}{\mid}
\optionalbracket{\interiorset}{S'}
\optionalbracket{\dirac}{\delta_x}
\newcommand{\generates}{\vdash}
\newcommand{\generatesNot}{\cancel{\vdash}}
\newcommand{\cond}[2]{#1(#2|C)}
\newcommand{\thick}{\;}

\newcommand{\disintegrates}{\div}
\newcommand{\functionOf}{\triangleleft}
\newcommand{\functionTo}{\triangleright}
\newcommand{\without}{\setminus}

\optionalbracket{\PA}{\text{Pa}}
\newcommand{\pa}{\text{pa}}
\optionalbracket{\range}{\text{range}}

\newcommand{\product}{\bigotimes}%

% definitions added by Magdalena
\newcommand{\FS}{\Omega} %\mathcal{F}} % factored space model
\newcommand{\FSModel}{\mathcal{M}} % factored space model
\newcommand{\Omg}{\Omega} % factors and their product
\newcommand{\omg}{\omega} % factors and their product
\newcommand{\U}{U} % projections/exogenous variables
\newcommand{\Obs}{{\textnormal{Obs}}} % observation space
\newcommand{\A}{\text{an}} % Ancestors
\newcommand{\Val}{\textnormal{Val}} % Value Space/Range
\newcommand{\Ainc}{\A_{\text{inc}}} % Inclusive Ancestors
\newcommand{\Vpriv}{{V}}%^{\textnormal{Priv}}}} % privileged variables
\newcommand{\Bayes}{\mathcal{B}}% bayes net
\newcommand{\Vars}{\text{Var}}%
\newcommand{\V}{\mathbf{V}}%
\newcommand{\Xb}{X}%\mathbf{X}}%
\newcommand{\pav}{{\pa(v)}}%
\newcommand{\xpav}{x_\pav}%
\newcommand{\xpavstar}{\xpav^*}%
\newcommand{\Xpav}{X_\pav}%
\newcommand{\vxpavstar}{{(v,\xpav^*)}}%
\newcommand{\vxpav}{{(v,\xpav)}}%
\newcommand{\Pomg}{P^\Omg}%
\newcommand{\iinI}{{i\in I}}%
\newcommand{\Ieq}{{I_{=}^x}}%
\newcommand{\Ineq}{{I_{\ne}^x}}%
\newcommand{\iinIeq}{{i\in \Ieq}}%
\newcommand{\iinIneq}{{i\in \Ineq}}%
\newcommand{\vinV}{{v\in V}}%
\newcommand{\merge}{\cdot}%
\newcommand{\timesOmg}{\bigtimes_\iinI\Omg_i}%
\newcommand{\FSG}{\FS}%
\newcommand{\x}{\alpha}%
\newcommand{\y}{\beta}%
\newcommand{\FSMG}{\FSModel^G}%
\newcommand{\Acal}{\mathcal{A}}%


\renewcommand{\epsilon}{\varepsilon}%


\makeatletter
\renewcommand{\paragraph}{\@ifstar\@paragraphbold\@paragraphbold}
\newcommand{\@paragraphbold}[1]{\vspace{4mm}\noindent\textbf{{#1}.}}
\makeatother


\optionalbracket{\cvar}{C}
\optionalbracket{\base}{\text{Base}}
\optionalbracket{\history}{\mathcal{H}}
\optionalbracket{\relevant}{\mathcal{R}}
\optionalbracket{\supp}{\mathrm{supp}}
\optionalbracket{\proj}{\pi}
\optionalbracket{\marg}{\textnormal{marg}}



\newcommand{\flatunion}{\union}
\newcommand{\flatunionbig}{\product}
\newcommand{\flattimes}{\times}
\newcommand{\flattimesbig}{\times}
\newcommand{\proddistr}{\product}
\newcommand{\condvar}[2]{#2}
\newcommand{\inrange}[2]{#1}
\newcommand{\casesinline}[2]{\begin{cases}#1\end{cases}}
\newcommand{\ms}{-measurable}
\newcommand{\NR}[1]{\NN_{#1}}
\newcommand{\num}[1]{#1}
\newcommand{\compose}{\circ}




\newcommand{\MS}{\mathcal{M}}
\newcommand{\emb}{\mathfrak{e}}
\newcommand{\xor}{\oplus}
\newcommand{\before}{\leq}
\newcommand{\strictlybefore}{<}
\newcommand{\strictlybeforeF}{\strictlybefore^\FS}
\newcommand{\beforeF}{\before^\FS}
\newcommand{\equivalent}{\simeq}
\newcommand{\Var}[1]{V_{#1}}
\newcommand{\model}{m}


\DeclareMathOperator{\historyC}{Cohistory}



\newtcolorbox{definitiontext}[1][]{
  colback=white,
  colframe=darkgray,
  boxrule=0.5mm,
  #1
}
